\chapter{Introducao}
\label{chap_introducao}

Este trabalho atende ao roteiro proposto para a pratica de Automatos de Wolfram: implementar um codigo executavel que permita selecionar a regra, definir estado inicial e controlar a densidade e o agrupamento dos sitios ocupados, produzindo como saida a evolucao espaco-temporal do sistema.

Nesta entrega, o fluxo adotado foi centrado no aplicativo/executavel da interface grafica. Nao foi utilizado notebook \texttt{.ipynb} para gerar os resultados apresentados.

As simulacoes foram desenvolvidas em Python e estruturadas para:
\begin{itemize}
\item executar regras selecionadas e regras classificadas por classe de comportamento;
\item gerar e salvar os graficos correspondentes;
\item disponibilizar uma interface grafica empacotada em executavel para uso sem depender da execucao manual do script Python.
\end{itemize}

No conjunto final de experimentos, foram utilizadas regras representativas das quatro classes de Wolfram (homogenea, periodica, caotica e complexa), com analise qualitativa por figuras e discussao baseada na literatura classica da area \cite{wolfram1983,wolfram2002,mathworldECA,wikipediaECA}.

O codigo-fonte desta pratica esta publicado no GitHub, no repositorio do projeto, junto com os artefatos de simulacao e documentacao \cite{githubProjetoWolfram}.
